\documentstyle[%


righttag%


,11pt%


,amscd%


,amssymb%


,russian%               remove in Eng.


,izvru%                 Set izven% in Eng.


]{amsart}





%%% NB !!!!!! If you want to compile this file, first remove


%%% a period (.) in line 522 of ANSART.STY (...\LATEX209)


%%% and then return it back, since it will affect ALL other files and


%%% environments


%%% remove \hspaces in eng.


% change baseline...


% attn to limits!


\newcommand{\sign}{\operatorname{sign}}


\newcommand{\im}{\operatorname{Im}}


\newcommand{\tts}[1]{}


%\renewcommand{\baselinestretch}{0.98}





\theoremstyle{plain}


\theorembodyfont{\it}


\newtheorem{assrt}{\hskip\parindent Утверждение}





%\hoffset=-15mm%                  For Eng.


\setcounter{page}{33}


\god{2000}


\nomer{10~(461)} %                        Remove in Eng.


%\nomer{Vol.\,44, No.\,10}%               Set in Eng.


%\str{\pageref{begin}--\pageref{end}}%  Set in Eng.


\udk{517.518; 517.982}





\author{\it Т.П.\,Лукашенко}


% NB:   ^^ remove in Eng.


\title{О свойствах обобщенных систем


разложения, \\ подобных ортогональным}


\thanks{Работа выполнена при финансовой поддержке Российского фонда


фундаментальных исследований (грант \No\,99-01-00354) и программы


поддержки ведущих научных школ РФ (грант \No\,00-15-96143).}





\date{}


%\pagestyle{myheadings}%         Set in Eng.


\pagestyle{plain} %              Remove in Eng.


\begin{document}


%\thispagestyle{plain}%          Set in Eng.


\maketitle


\label{begin}








\section{Введение}





Известно, что ортогональные системы и разложения по ним широко


применяются в математике и ее приложениях, что объясняется рядом


замечательных свойств таких систем. В [1] и


[2] было дано определение более общих систем разложения, подобных


ортогональным (более кратко, ортоподобных).





\begin{defn}\hspace{-1.8mm}{\bf.}


Пусть $H$ --- гильбертово пространство над


полем $\Bbb R$ или $\Bbb C$,


а $\Omega$ --- пространство со счетно-аддитивной (соответственно


действительной или комплексной) мерой $\mu$ ([3],


с.\,109--116). Систему $\{e^\omega\}_{\omega\in\Omega}\subset H$


будем называть {\it системой в $H$ с индексами из $\Omega$}.


\end{defn}





\begin{defn}\hspace{-2mm}{\bf.}


Систему $\{e^\omega\}_{\omega\in\Omega}$


в $H$ с индексами из $\Omega$ будем называть {\it ортоподобной


$($подобной ортогональной$)$ системой разложения} в $H$, если любой


элемент $y\in H$ можно представить в виде


\begin{equation}


y=\int_\Omega\widehat y_\omega e^\omega\,d\mu(\omega),


\end{equation}


где $\wh y_\omega=(y,e^\omega)$, интеграл понимается как


собственный или несобственный интеграл Лебега от функции со


значениями в $H$, причем в последнем случае имеется такое исчерпывание


$\{\Omega_k\}^\infty_{k=1}$ пространства $\Omega$ (все $\Omega_k$


измеримы, $\Omega_k\subset\Omega_{k+1}$ для $k\in\Bbb N$ и


$\bigcup^\infty_{k=1}\Omega_k =\Omega$), быть может, зависящее


от $y$ и называемое подходящим для $y$, что функция $\wh


y_\omega e^\omega$ интегрируема по Лебегу на $\Omega_k$ и


\begin{equation}


y=\int_\Omega\wh y_\omega e^\omega\,d\mu(\omega)=


\lim_{k\to\infty}(L)\int_{\Omega_k}\wh y_\omega e^\omega


\,d\mu(\omega).


\end{equation}


\end{defn}





Ортонормированные системы являются ортоподобными системами


разложения в замыкании своей линейной оболочки; при этом для любой


точки $\omega\in\Omega\quad\mu(\omega)\equiv 1$ и для счетных систем


в качестве $\Omega$ можно брать $\Bbb N$ --- множество натуральных


чисел.


Для ортоподобных систем разложения имеет место





\begin{assrt} {\em (аналог равенства


Парсеваля) ([2], [4], [5])}. Если


$\{e^\omega\}_{\omega\in\Omega}$ --- ортоподобная система


разложения в $H$, то для любого элемента $y\in H$


$$


\|y\|^2 =\int_\Omega |\wh y_\omega|^2\,d\mu(\omega


)=\int_\Omega |(y,e^\omega )|^2\,d\mu(\omega),


$$


где в соответствии с определением интеграл понимается как собственный


интеграл Лебега или как несобственный интеграл Лебега по подходящему


для элемента $y$ исчерпыванию$;$ для любых элементов $y,z\in H$


$$


(y,z)=\int_\Omega \wh y_\omega\ov{\wh z_\omega}\,


d\mu(\omega)=\int_\Omega (y,e^\omega)\ov{(z,e^\omega)}\,


d\mu(\omega),


$$


где также интеграл понимается как собственный интеграл Лебега или


как несобственный интеграл Лебега по подходящему для $y$ или $z$


исчерпыванию. Если мера $\mu$ неотрицательна, то в обоих случаях


интеграл может рас\-смат\-ри\-вать\-ся как собственный интеграл


Лебега.


\end{assrt}





Из утверждения 1 выводится





\begin{assrt} {\em [2] (о мнимой части меры)}.


Если $\{e^\omega\}_{\omega\in\Omega}$ --- ортоподобная система


разложения в $H$, то $\int\limits_\Omega\wh y_\omega e^\omega\,d


\im \mu(\omega)=0$ для любого элемента $y\in H$, поэтому


мнимую часть меры $\mu$ можно отбросить и без ограничения общности


считать $\mu$ действительной мерой.


\end{assrt}





\begin{defn}\hspace{-1.8mm}{\bf.}


Систему элементов $\{e^\omega\}_{\omega


\in\Omega}$ из $H$ с индексами из $\Omega$ будем называть {\it


неотрицательной (с неотрицательной мерой)}, если $\Omega$ ---


пространство со счетно-аддитивной неотрицательной мерой $\mu$.


\end{defn}





Всюду ниже функции на $\Omega$ со значениями в $\Bbb R$ или


$\Bbb C$ всегда принимают значения в том поле, над которым


рассматривается $H$.





В случае неотрицательной меры $\mu$ верны следующие


доказанные в [4] и [5] утверждения.





\begin{assrt} {\em (об экстремальном свойстве коэффициентов


разложения)}. Пусть $\{e^\omega\}_{\omega\in\Omega}$ ---


неотрицательная ортоподобная система, а $c(\omega)$ --- функция на


$\Omega$ со значениями в $\Bbb R$ или $\Bbb C$~и


\begin{equation}


y=\int_\Omega c(\omega)e^\omega\,d\mu(\omega),


\end{equation}


где интеграл понимается как собственный или несобственный интеграл


Лебега от функции со значениями в $H$, причем в последнем случае


существует такое исчерпывание $\{\Lambda_k\}^\infty_{k=1}$


пространства $\Omega$ $($все $\Lambda_k$ измеримы, $\Lambda_k\subset


\Lambda_{k+1}$ для $k\in \Bbb N$ и $\bigcup^\infty_{k=1}


\Lambda_k=\Omega)$, что функция $c(\omega)e^\omega$ интегрируема


по Лебегу на $\Lambda_k$ и


$$


\int_\Omega c(\omega)e^\omega\,d\mu(\omega)=\lim_{k\to


\infty}(L)\int_{\Lambda_k}c(\omega)e^\omega\,d\mu(\omega).


$$


Тогда выполняется следующий аналог равенства Парсеваля$:$ для любого


элемента $z\in H$


\begin{equation}


(y,z)=\int_\Omega с(\omega)\ov{\wh z_\omega}\,d\mu(\omega


)=\int_\Omega с(\omega)\ov{(z,e^\omega)}\,d\mu(\omega),


%\tag{4}


\end{equation}


где интеграл также понимается как собственный или соответственно


несобственный интеграл Лебега --- предел интегралов по вышеуказанному


исчерпыванию $\{\Lambda_k\}^\infty_{k=1}$ пространства $\Omega$. Если


$\int\limits_\Omega |c(\omega)|^2\,d\mu(\omega)<\infty$, то в любом случае


интеграл может рассматриваться как собственный интеграл Лебега.


Полагая $z=y$, имеем равенство $\|y\|^2=\int\limits_\Omega c(\omega


)\ov{\wh y_\omega}\,d\mu(\omega)$, в котором заменить


$\wh y_\omega$ на $c(\omega)$, вообще говоря, нельзя$:$


$$


\|y\|^2\leqslant(L)\int_\Omega|c(\omega)|^2\,d\mu(\omega),


$$


причем равенство имеет место лишь в случае $c(\omega)=\wh


y_\omega$ почти всюду на $\Omega$.


\end{assrt}





{\bf Замечание.} Для выполнения только аналога равенства


Парсеваля (4) неотрицательность меры $\mu$ не нужна.





\begin{assrt} {\em (аналог теоремы Рисса--Фишера)}. Пусть


$c(\omega)$ --- функция на $\Omega$ со значениями в $\Bbb R$ или


$\Bbb C$ и %%%%%% limits


$\int_\Omega|c(\omega)|^2\,d\mu(\omega)<\infty$.


Пусть ${\{\Lambda_k}\}^\infty_{k=1}$ --- такая последовательность


измеримых подмножеств $\Omega$, что $\varliminf_{k\to\infty} %%%% limits


{\Lambda_k}=\Omega$ {\em ([3], с.\,41)} и функция $c(\omega


)e^\omega$ интегрируема по Лебегу на $\Lambda_k$. Тогда в $H$


существует


$$


\lim_{k\to\infty}(L)\int_{\Lambda_k}c(\omega)e^\omega\,d\mu(\omega).


$$


\end{assrt}





\begin{assrt} {\em (оценка точности приближения)}. Пусть


$c(\omega)$ --- функция на $\Omega$ со значениями в $\Bbb R$ или


$\Bbb C$ и $\int\limits_\Omega|c(\omega)|^2\,d\mu(\omega)<\infty$,


$E$ --- измеримое подмножество $\Omega$. Тогда при условии измеримости


функции $c(\omega)e^\omega$ на $\Omega$ существуют собственные


или несобственные интегралы Лебега от функции $c(\omega)e^\omega$ по


$\Omega$ и по $E$ и


$$


\bigg\| \int_\Omega c(\omega)e^\omega\,d\mu(\omega)-


\int_E c(\omega)e^\omega\,d\mu(\omega)\bigg\|^2\leqslant


\int_{\Omega\setminus E}|c(\omega)|^2\,d\mu(\omega).


$$


\end{assrt}





{\bf Замечание.} Если при условиях утверждения 5\;


$\int\limits_\Omega c(\omega)e^\omega\,d\mu(\omega)$ обозначить


через $y$, то оценка точности приближения будет выглядеть как


$$


\bigg\|y-\int_E c(\omega)e^\omega\,d\mu(\omega)\bigg\|^2


\leqslant\int_{\Omega\setminus E}|c(\omega)|^2\,d\mu(\omega).


$$


В частности, для любого элемента $y\in H$ и любого измеримого


подмножества $E$ из $\Omega$ существует собственный или


несобственный интеграл Лебега


$\int\limits_E\wh y_\omega e^\omega\,d\mu(\omega)$ и


$$


\bigg\| y-\int_E\wh y_\omega e^\omega\,d\mu(\omega)\bigg


\|^2\leqslant\int_{\Omega\setminus E}|\wh y_\omega|^2\,


d\mu(\omega).


$$





Это замечание непосредственно вытекает из утверждения 4 и


равенства (1). Последние оценки являются аналогами известной в


теории рядов Фурье оценки точности приближения --- тождества Бесселя.





Опубликованные [2], [4], [5] утверждения показывают, что системы


разложения, подобные ортогональным, обладают почти всеми свойствами


ортогональных систем. В случае, когда $H$ --- пространство Лебега


$L^2$, они обладают многими замечательными свойствами ортогональных


систем функций. Однако задающая преобразование Фурье система $\{\text


{exp}(2\pi i\omega x)\}_{\omega\in\Bbb R}$, определяющая


преобразование Гильберта система $\big\{\frac1{\pi(\omega-x)}


\big\}_{\omega\in\Bbb R}$ и некоторые системы разложения из [6] не


являются системами разложения, подобными ортогональным, хотя обладают


многими свойствами ортогональных систем. Поэтому естественным является


введение и изучение более общего понятия --- обобщенной системы


разложения, подобной ортогональной.





\section{Обобщенные ортоподобные системы разложения}





\begin{defn}\hspace{-1.8mm}{\bf.}


Пусть $H$ --- гильбертово пространство над


полем $\Bbb R$ или $\Bbb C$, а $\Omega$ --- пространство со


счетно-аддитивной (соответственно действительной или комплексной) мерой


$\mu$. Пусть $\{H_n\}_{n=1}^\infty$ --- система замкнутых


расширяющихся ($H_n\subset H_{n+1}$ для $n\in\Bbb N$)


подпространств в $H$, объединение которых всюду плотно в $H$, а


$\{\boldsymbol e^\omega\}_{\omega\in\Omega}$ --- такая система, каждый


элемент которой $\boldsymbol e^\omega$ является последовательностью


$\{e^\omega_n\}_{n=1}^\infty$ элементов $H$ и $e^\omega_n$ ---


ортогональная проекция $e^\omega_{n+1}$ на $H_n$. Тогда систему


$\{\boldsymbol e^\omega\}_{\omega\in\Omega}$ будем называть {\it


обобщенной системой в $H$ $($с системой подпространств


$\{H_n\}_{n=1}^\infty$ и индексами из $\Omega)$}.


\end{defn}





Из определения следует, что для любых $m,n\in\Bbb N$, $m>n$, и


любых $\omega\in\Omega$ элемент $e^\omega_n$ --- ортогональная


проекция $e^\omega_m$ на $H_n$.





\begin{defn}\hspace{-1.8mm}{\bf.}


Обобщенная система $\{\boldsymbol e^\omega


\}_{\omega\in\Omega}$ будет называться {\it обобщенной ортоподобной


$($подобной ортогональной$)$ системой разложения в $H$}, если любой


элемент $y\in H_n$ представляется в виде


\begin{equation}


y=\int_\Omega\wh y^n_\omega e^\omega_n\,d\mu(\omega),


\end{equation}


где $\wh y^n_\omega=(y,e^\omega_n)$, интеграл понимается как


собственный или несобственный интеграл Лебега от функции со


значениями в $H$ (точно так же, как в определении 2).


\end{defn}





\begin{defn}\hspace{-1.8mm}{\bf.}


Обобщенную систему $\{\boldsymbol e^\omega


\}_{\omega\in\Omega}$ будем называть {\it неотрицательной $($с


неотрицательной мерой$)$}, если $\Omega$ --- пространство со


счетно-аддитивной неотрицательной мерой~$\mu$.


\end{defn}





\begin{note} (об интегрировании). Используемое в


определениях 2 и 4 понятие интеграла является, по мнению автора,


наиболее пригодным для излагаемой далее теории, хотя попытки


использования других интегралов также могут быть интересны.


\end{note}





\begin{note} (о числовом множителе).


Разумеется, в интегралы


в формулах (1)--(5) можно добавить действительнозначный или


комплекснозначный множитель $a(\omega)$, но от него легко избавиться


переходом к системе $\{\boldsymbol s^\omega\}=\{\sqrt{|a(\omega


)|}\boldsymbol e^\omega\}$ и мере $\eta(E)=(L)\int\limits_E\frac


{a(\omega)}{|a(\omega)|}\,d\mu(\omega)$ (здесь считаем $\frac00=0$),


$$


(L)\int_{\Omega_{(k)}}a(\omega)(y,e^\omega_n)e^\omega_n\,


d\mu(\omega)=(L)\int_{\Omega_{(k)}}(y,s^\omega_n)s^\omega_n\,


d\eta(\omega),


$$


где $\Omega_{(k)}$ --- $\Omega$ или $\Omega_k$ ([3],


с.\,198). Существуют и другие переходы, также избавляющие от множителя


$a(\omega)$.


\end{note}





\begin{note} (об ортоподобных системах).


Если в обобщенной


ортоподобной системе разложения все элементы $e^\omega\equiv


\{e^\omega_n\}_{n=1}^\infty$ постоянные (соответственно все $H_n=H$),


то получается определение ортоподобной системы разложения, в


частности, охватывающее полные ортонормированные системы.


\end{note}





\begin{note} (о преобразовании Фурье).


Если взять $H=L^2


(\Bbb R)$, $\Omega=\Bbb R$ c $\mu$ --- обычной мерой Лебега,


$H_n=\{f\in L^2(\Bbb R): f\equiv 0\text{ на }\Bbb


R\setminus[-n,n]\}$, $e^\omega_n(x)=\exp(2\pi i\omega


x)\chi_{[-n,n]}(x)$, где $\chi_{[-n,n]}(x)$ --- характеристическая


функция отрезка $[-n,n]$, то в


качестве обобщенной ортоподобной неотрицательной системы получится


система функций $\{\exp(2\pi i\omega x)\}_{\omega\in\Bbb R}$,


определяющая преобразование Фурье $\cal F(f)(\omega)\equiv\wh


f(\omega)=\int\limits_{\Bbb R}\exp(-2\pi i\omega x)f(x)\,dx$.


\end{note}





\begin{note} (о преобразовании Гильберта).


Если взять


$H=L^2(\Bbb R)$, $\Omega=\Bbb R$ c $\mu$ --- обычной мерой


Лебега, $H_n=\{f\in L^2(\Bbb R):\cal F(f)\equiv 0$


на $\Bbb R\setminus[-n,n]\}$,


$$


e^\omega_n(x)=\cal F(-i\exp(2\pi i\omega t)


\sign t\,\chi_{[-n,n]}(t))(x)=\frac{1-\cos 2\pi


n(\omega-x)}{\pi(\omega-x)},


$$


то в качестве обобщенной ортоподобной


неотрицательной системы получится система функций


$\big\{\frac1{\pi(\omega-x)}\big\}_{\omega\in\Bbb R}$,


определяющая преобразование Гильберта $\cal H(f)(\omega)\equiv


\wt f(\omega)=\text{v.p.}\int\limits_{\Bbb R}\frac{f(x)}{\pi(\omega-x)}


\,dx\equiv\lim\limits_{\varepsilon\to+0}


\int\limits_{|x-\omega|\geqslant\varepsilon}\frac{f(x)}


{\pi(\omega-x)}\,dx$ ([7], с.\,243).


\end{note}





Из приведенных определений 1 и 2 следует, что для любого $n\in


\Bbb N$ система $\{e^\omega_n\}_{\omega\in\Omega}$ является


ортоподобной системой разложения в $H_n$, неотрицательной, если


обобщенная система разложения $\{\boldsymbol e^\omega\}_{\omega\in


\Omega}$ неотрицательна.





Рассматривая обобщенные системы разложения, подобные


ортогональным, будем в дальнейшем обозначать через $P_n$ оператор


ортогонального проектирования из $H$ в $H_n$. Известно ([8],


с.\,139--145), что для любого элемента $y\in H$


элемент $P_n(y)$ является ближайшим к $y$ элементом $H_n$. Так как


в определении 4 замыкание ${\bigcup_{n=1}^\infty H_n}$ равно


$H$, то для любого $y\in H$ верно, что $y=\lim\limits_{n\to\infty}P_n(y)$.





\begin{thm}\hspace{-1.8mm}{\bf.}


Если $\{\boldsymbol e^\omega\}_{\omega\in


\Omega}$ --- обобщенная ортоподобная неотрицательная система


разложения в $H$, то для любого элемента $y\in H$ существует


единственная с точностью до эквивалентности $($совпадения почти всюду$)$


функция $\wh y_\omega$ на $\Omega$ такая, что последовательность


функций $\wh y^n_\omega=(y,e^\omega_n)$ сходится к ней в смысле


$$


\lim_{n\to\infty}\int_\Omega|\wh y_\omega-\wh


y_\omega^n|^2\,d\mu(\omega)=0.


$$


Выполняется аналог равенства Парсеваля--Планшереля$:$ для любого


элемента $y\in H$


$$


\| y\|^2 =(L)\int_\Omega|\wh y_\omega|^2\,d\mu(\omega);


$$


для любых элементов $y,z\in H$


$$


(y,z)=(L)\int_\Omega\wh y_\omega\ov{\wh


z_\omega}\,d\mu (\omega ),


$$


где $(L)$ перед знаком интеграла означает, что


интеграл понимается как собственный интеграл Лебега.


\end{thm}





{\bf Доказательство.}


Для любых $n,m\in\Bbb N$, $n\geqslant m$, в силу аналога равенства


Парсеваля для систем разложения, подобных ортогональным ([2] или


[4], или [5]),


\begin{multline*}


\|P_n(y)-P_m(y)\|^2=(L)\int_\Omega|(


P_n(y),e^\omega_n)-(P_m(y),e^\omega_n)|^2\,d\mu(\omega)= \\


%


=(L)\int_\Omega|(y,e^\omega_n)-(P_m(y),e^\omega_m)|^2


\,d\mu(\omega)=(L)\int_\Omega|\wh y^n_\omega


-\wh y^m_\omega|^2\,d\mu(\omega).


\end{multline*}





Так как замыкание $\bigcup_{n=1}^\infty H_n$ совпадает с $H$,


то $P_n(y)@>>{n\to\infty}>y$. Поэтому для любого


$\varepsilon>0$ найдется такое натуральное $N$, что для любых $n,m>N$


\begin{equation}


(L)\int_\Omega|\wh y^n_\omega-\wh y^m_\omega


|^2\,d\mu(\omega)<\varepsilon.


\end{equation}





Выберем такую строго возрастающую последовательность натуральных


чисел $\{n_k\}$, что $(L)\int\limits_\Omega|\wh y_\omega^{n_{k+1}}


-\wh y_\omega^{n_k}|^2\,d\mu(\omega )<8^{-k}$, $k\in\Bbb


N$. Тогда $\mu\{\omega\in\Omega: |\wh y^{n_{k+1}}_\omega


-\wh y^{n_k}_\omega|>2^{-k}\}<2^{-k}$ и, значит,


последовательность $\wh y^{n_k}_\omega$ сходится почти


всюду на $\Omega$ к функции, которую обозначим $\wh y_\omega$.


Неотрицательные интегрируемые по Лебегу функции $|\wh


y^n_\omega-\wh y^m_\omega|^2$ при $n\to\infty$ сходятся


почти всюду на $\Omega$ к функции $|\wh y_\omega-\wh


y^m_\omega|^2$. Пользуясь леммой Фату ([3],


с.\,170), из (6) получаем, что при $m>n$ выполняется неравенство


$(L)\int\limits_\Omega|\wh y_\omega-\wh y^n_\omega|^2\,


d\mu(\omega)<\varepsilon$. Следовательно,


$$


\lim_{m\to\infty}\int_\Omega|\wh y_\omega-\wh 


y_\omega^m|^2\,d\mu(\omega)=0.


$$


По аналогу равенства Парсеваля для систем разложения, подобных


ортогональным, для любых элементов $y,z\in H$ и $m\in\Bbb N$


$$


%\begin{equation*}


\begin{CD}


(P_m(y),P_m(z))@=(L)\displaystyle\int_\Omega\wh y^m_\omega


\ov{\wh z^m_\omega}\,d\mu (\omega ) \\


@VV{m\to\infty}V @VV{m\to\infty}V \\


(y,z) @.(L)\displaystyle\int_\Omega\wh y_\omega


\ov{\wh z_\omega}\,d\mu (\omega )


\end{CD}


$$


%\end{equation*}


и, следовательно, для любых элементов $y,z\in H$ выполняется аналог


равенства Парсеваля--Планшереля


$$


(y,z)=(L)\int_\Omega\wh y_\omega\ov{\wh


z_\omega}\,d\mu (\omega).\quad\square


$$





Доказанный аналог равенства Парсеваля--Планшереля существенно


используется в нижеследующих результатах.





\begin{thm}\hspace{-1.8mm}{\bf.}


Пусть $\{\boldsymbol e^\omega\}_{\omega\in


\Omega}$ --- неотрицательная обобщенная ортоподобная система


разложения в $H$, а $c(\omega)$ --- функция со значениями в $\Bbb


R$ или в $\Bbb C$, $|c(\omega)|^2$ интегрируема по Лебегу на


$\Omega$ и


$$


y=\lim_{n\to\infty}\int_\Omega c(\omega)e^\omega_n\,d\mu(\omega),


$$


где интегралы понимаются как собственные или несобственные интегралы


Лебега от функции со значениями в $H$. Тогда для любого $z\in H$


$$


(y,z)=(L)\int_\Omega c(\omega)\ov{\wh z_\omega}\,


d\mu(\omega).


$$


\end{thm}





{\bf Доказательство.} Так как


скалярное произведение является при фиксированном втором аргументе


непрерывным линейным оператором от первого аргумента, то согласно


([3], с.\,128)


$$


(y,z)=\lim_{n\to\infty}\int_\Omega c(\omega)(e^\omega_n,z)\,


d\mu(\omega).


$$


Так как для любого $\varepsilon>0$ выполняется неравенство


$|c(\omega)(e^\omega_n,z)-c(\omega)\wh z_\omega|\leqslant


\frac12(\varepsilon^2|c(\omega)|^2+\varepsilon^{-2}|(e^\omega_n


,z)-\wh z_\omega|^2)$, то из теоремы 1 следует, что


$$


(y,z)=\int_\Omega c(\omega)\wh z_\omega\,d\mu(\omega).\quad\square


$$





{\bf Следствие} (об экстремальном свойстве коэффициентов


разложения). При условии теоремы~2


$$


\|y\|^2\leqslant(L)\int_\Omega|c(\omega)|^2\,d\mu(\omega),


$$


причем равенство имеет место тогда и только тогда, когда $c(\omega


)=\wh y_\omega$ почти всюду на $\Omega$.


\medskip





Действительно, полагая в теореме 2\; $z=y$, имеем $\|y\|^2


=(L)\int\limits_\Omega c(\omega)\ov{\wh y_\omega}\,d\mu(\omega)$.


Используя неравенство $|c(\omega)\ov{\wh y_\omega}


|\leqslant\frac12(|c(\omega)|^2+|\wh y_\omega|^2


)$ и теорему 1, получаем $\|y\|^2\leqslant\frac12


(L)\int\limits_\Omega |c(\omega)|^2\,d\mu(\omega)+\frac12\|y\|^2$, причем


равенство имеет место тогда и только тогда, когда почти всюду на


$\Omega$ выполняются равенство $|c(\omega)|=|\wh y_\omega|$


и $c(\omega)\ov {\wh y_\omega}\geqslant0$, т.\,е. тогда и


только тогда, когда $c(\omega)=\wh y_\omega$ почти всюду на~$\Omega$.





\begin{lem}\hspace{-1.8mm}{\bf.}


Пусть в гильбертовом пространстве $H$ имеется


система замкнутых подпространств $\{H_n\}_{n=1}^\infty$, $H_1\subset


H_2\subset H_3\subset\dotsb$, замыкание $\bigcup_{n=1}^\infty H_n$


совпадает с $H$, и $P_n$ --- оператор ортогонального проектирования из


$H$ в $H_n$. Если $y^n$ --- такая ограниченная последовательность в $H$,


что при $n<m$\; $y^n=P_n(y^m)$, то $y^n$ --- сходящаяся в $H$


последовательность с пределом $y$ и $y^n=P_n(y)$ для всех $n\in\Bbb


N$.


\end{lem}





\begin{pf}


При $m<n$\; $\|y^m\|


=\|P_m(y^n)\|\leqslant\|y^n\|$, значит, $\|y^k\|$ --- ограниченная


неубывающая последовательность, пусть $s$ --- ее предел. Возьмем


произвольное $\varepsilon>0$ и найдем такое $N$, что при


$k>N$\; $\|y^k\|^2>s^2-\varepsilon^2$.


Тогда при $n\geqslant m>N$ имеем $\|y^n


-y^m\|^2=\|y^n-P_m(y^n)\|^2=\|y^n\|^2-\|y^m\|^2<\varepsilon^2$.


Последовательность $y^k$ фундаментальна и, следовательно, сходится.


При $n<m$ имеем равенство $y^n=P_n(y^m)$. При $m\to\infty$


в силу непрерывности оператора ортогонального проектирования


получаем $y^n=P_n(y)$ для всех $n\in\Bbb N$.


\end{pf}





\begin{thm}\hspace{-1.8mm}{\bf.}


Пусть $\{\boldsymbol e^\omega\}_{\omega\in


\Omega}$ --- обобщенная ортоподобная неотрицательная система


разложения в $H$, тогда для любого $y\in H$ верно равенство


\begin{equation}


y=\lim_{n\to\infty}\int_\Omega \wh y^n_\omega


e_n^\omega\,d\mu(\omega).


\end{equation}


\end{thm}





\begin{pf}


Так как $(y,e_n^\omega)=(P_n(y),e^n_\omega


)$, по определению 5\; $P_n(y)=\int\limits_\Omega \wh y^n_\omega


e_n^\omega\,d\mu(\omega)$, а по лемме 1\; $y=\lim\limits_{n\to\infty}


P_n(y)$, то утверждение теоремы верно.


\end{pf}





\begin{thm}\hspace{-1.8mm}{\bf.}


Пусть $\{\boldsymbol e^\omega\}_{\omega\in


\Omega}$ --- обобщенная ортоподобная неотрицательная система


разложения в $H$, $c(\omega)$ --- функция со значениями в $\Bbb R$


или $\Bbb C$, $|c(\omega)|^2$ интегрируема по Лебегу на $\Omega$,


и для каждого $n\in\Bbb N$\; $\{\Lambda^n_k\}^\infty_{k=1}$ ---


такая последовательность измеримых подмножеств $\Omega$, что еe нижний


предел {\em ([3],


с.\,141)} $\underset{k\to \infty}{\varliminf}{\Lambda^n_k}=\Omega$ и функция


$c(\omega)e_n^\omega$ интегрируема по Лебегу на $\Lambda^n_k$. Тогда в


$H$ существует


\begin{equation}


\lim_{n\to\infty}\lim_{k\to\infty}(L)\int_{\Lambda_k}c(\omega


)e_n^\omega\,d\mu(\omega).


\end{equation}


\end{thm}





\begin{pf}


По теореме 3 из [4] для каждого $n\in\Bbb N$


существует $y^n=(L)\int\limits_{\Lambda^n_k}c(\omega)e_n^\omega\,


d\mu(\omega)$. Так как при $m\leqslant n$ имеет место равенство


$P_m(y^n)=(L)\int\limits_{\Lambda^n_k}c(\omega)P_m(e_n^\omega)\,d\mu(\omega)


=(L)\int\limits_{\Lambda^n_k}c(\omega)e_m^\omega\,d\mu(\omega)$ ([3],


с.128), а по теореме 2 из [4] для любого $n\in\Bbb N$\;


$\|y^n\|^2\leqslant\int\limits_\Omega|c(\omega)|^2\,d\mu(\omega)<\infty$,


то по лемме~1 последовательность $y^n$ сходится.


\end{pf}





Доказанная теорема является аналогом известной для


ортонормированных систем теоремы Рисса--Фишера ([3], с.\,414).





\begin{cor}\hspace{-1.8mm}{\bf.}


Если $\{\boldsymbol e^\omega\}_{\omega\in


\Omega}$ --- обобщенная ортоподобная неотрицательная система


разложения в $H$, то для любого $y\in H$ верно равенство


\begin{equation}


y=\lim_{n\to\infty}\int_\Omega \wh y_\omega


e_n^\omega\,d\mu(\omega).


\end{equation}


\end{cor}





Действительно, $P_n(y)=\int\limits_\Omega \wh y^n_\omega


e_n^\omega\,d\mu(\omega)$, следовательно, по теореме 1


$$


\bigg\|P_n(y)-\int_\Omega\wh y_\omega e_n^\omega\,d\mu


(\omega)\bigg\|^2=\int_\Omega |\wh y_\omega^n


-\wh y_\omega|^2\,d\mu(\omega)@>>{n\to\infty}>0


$$


и по теореме 3 следствие верно.


\medskip





{\bf Замечание.} Если в формуле (7) рассматривать внутренний


предел как несобственный интеграл Лебега, то при условии теоремы 3


в $H$ существует $\lim\limits_{n\to\infty}\int\limits_\Omega c(\omega


)e_n^\omega\,d\mu(\omega)$, который естественно обозначить короче.





\begin{defn}\hspace{-1.8mm}{\bf.}


Будем в дальнейшем $\lim_{n\to\infty}


\int_\Omega c(\omega)e_n^\omega\,d\mu(\omega)$, в котором


интеграл понимается как несобственный интеграл Лебега, обозначать


$$


\int_\Omega c(\omega)\boldsymbol e^\omega\,d\mu(\omega)


$$


и называть интегралом от обобщенной ортоподобной неотрицательной


системы разложения $\{\boldsymbol e^\omega\}_{\omega\in\Omega}$ с


коэффициентами $c(\omega)$.


Тем самым сможем короче записывать равенства (7)--(9).


\end{defn}





\begin{cor} (оценка точности приближения). Пусть


$c(\omega)$ --- функция со значениями в $\Bbb R$ или $\Bbb C$,


$|c(\omega)|^2$ интегрируема по Лебегу на $\Omega$, функция $c(\omega)


\boldsymbol e^\omega$ измерима на $\Omega$ и $E$ --- измеримое


подмножество $\Omega$. Тогда существуют собственные или несобственнные


(как в примечании выше) интегралы Лебега от функции


$c(\omega)\boldsymbol e^\omega$ по $\Omega$ и по $E$ и


$$


\bigg\| \int_\Omega c(\omega)\boldsymbol e^\omega\,d\mu(\omega)-


\int_E c(\omega)\boldsymbol e^\omega\,d\mu(\omega)\bigg\|^2\leqslant


\int_{\Omega\setminus E}|c(\omega)|^2\,d\mu(\omega).


$$


\end{cor}





{\bf Замечание.} Если при условиях следствия


$\int\limits_\Omega c(\omega)\boldsymbol e^\omega\,d\mu(\omega)$ обозначить


через $y$, то оценка точности приближения будет выглядеть как


$$


\bigg\|y-\int_E c(\omega)\boldsymbol e^\omega d\mu(\omega)\bigg\|^2\leqslant


\int_{\Omega\setminus E}|c(\omega)|^2d\mu(\omega).


$$


В частности, для любого элемента $y\in H$ и любого $E$ ---


измеримого подмножества $\Omega$


$$


\bigg\| y-\int_E\wh y_\omega\boldsymbol e^\omega\,d\mu(\omega


)\bigg\|^2\leqslant\int_{\Omega\setminus E}|\wh y_\omega|^2\,


d\mu(\omega).


$$





Это замечание непосредственно вытекает из утверждения 4 и


равенства (1). Последние оценки являются аналогами известной в


теории рядов Фурье оценки точности приближения --- тождества Бесселя.





\section{Вопросы измеримости ортоподобных систем.\protect\\


Пространство коэффициентов разложений}





В ряде мест предыдущего параграфа автор фактически повторяет


известные в теории пространств Лебега $L^p$ доказательства и


рассуждения, т.\,к. прямая ссылка недопустима из-за возможной


неизмеримости функций $\wh y_\omega^n=(y,e^\omega_n)$. В самом


деле, если $\{\boldsymbol e^\omega\}_{\omega\in\Omega}\subset H$ ---


обобщенная ортоподобная система разложения в $H$, а $\theta(\omega)$


--- функция на $\Omega$ со значениями в $\Bbb R$ или $\Bbb C$,


$|\theta(\omega)|\equiv 1$ на $\Omega$, то для любого элемента $y\in


H_n$ и любых $\omega\in\Omega$ и $n\in\Bbb N$ верно равенство


$(y,\theta(\omega)e^\omega_n)\theta(\omega)e^\omega_n=(y,e^\omega_n


)e^\omega_n$ и, значит, система $\{\theta(\omega)\boldsymbol e^\omega


\}_{\omega\in\Omega}$ --- также ортоподобная система разложения в $H$,


но функции $(y,\theta(\omega)e^\omega_n)=\ov{\theta(\omega)}


(y,e^\omega_n)$ при этом могут оказаться неизмеримыми. Вместе с


функциями $\wh y_\omega^n$ неизмеримой может быть и функция


$\wh y_\omega$. В связи с этим введем





\begin{defn}\hspace{-1.8mm}{\bf.}


Неотрицательную обобщенную ортоподобную


систему $\{\boldsymbol e^\omega\}_{\omega\in\Omega}$ будем


называть {\it измеримой}, если все функции $\wh y^n_\omega


=(y,e^\omega_n)$ измеримы.


\end{defn}





В случае счетно-конечного пространства с мерой $\Omega$ любую


неотрицательную обобщенную ортоподобную систему $\{e^\omega_n


\}_{n\in\Bbb N,\,\omega\in\Omega}$ можно сделать измеримой


умножением на функцию $\theta(\omega)$, по модулю равную единице.


Для доказательства этого понадобится





\begin{lem} {\em ([5])}.


Если пространство $\Omega$ со счетно-аддитивной


неотрицательной мерой $\mu$ счетно-конечно, т.\,е. является не более чем


счетным объединением множеств конечной меры, а $\Xi$ --- множество


функций на $\Omega$ со значениями в $\Bbb R$ или $\Bbb C$,


причем для любых двух функций $\varphi(\omega),\psi(\omega)\in\Xi$


произведение $\varphi(\omega)\ov {\psi(\omega)}$ измеримо на


$\Omega$, то существует такая функция $\theta(\omega)$ на $\Omega$


со значениями соответственно в $\Bbb R$ или $\Bbb C$, $|\theta


(\omega)|=1$ на $\Omega$, что для любой функции $\varphi(\omega)\in


\Xi$ функция $\ov {\theta(\omega)}\varphi(\omega)$ измерима на


$\Omega$.


\end{lem}





\begin{thm} {\em (об измеримости системы)}.


Если


$\{\boldsymbol e^\omega\}_{\omega\in\Omega}\equiv\{e^\omega_n\}_{n\in


\Bbb N, \omega\in\Omega}$ --- неотрицательная обобщенная


ортоподобная система разложения в $H$, а пространство с мерой $\Omega$


счетно-конечно, то существует такая функция $\theta(\omega)$ на


$\Omega$ со значениями в $\Bbb R$ или $\Bbb C$, $|\theta(\omega


)|=1$ на $\Omega$, что $\{\theta(\omega)e^\omega_n\}_{n\in\Bbb


N,\omega\in\Omega}$ --- измеримая неотрицательная обобщенная


ортоподобная система разложения в $H$.


\end{thm}





\begin{pf}


В силу определения 5


для любого $y\in H$ и любого $n\in\Bbb N$ имеет место равенство


$\wh y^n_\omega\equiv(y, e_n^\omega)=(P_n(y),


P_n(e_k^\omega))=(P_n(y),e_k^\omega)$, где $k>n$,


и, следовательно, по утверждению 1 для любых элементов $y,z\in H$ и


для любых $n,m\in\Bbb N\quad\wh y^n_\omega\ov{\wh


z^m_\omega}$ --- интегрируемая по Лебегу на $\Omega$ функция,


$$


(L)\int_\Omega \wh y^n_\omega\ov{\wh


z^m_\omega }\,d\mu (\omega )=(L)\int_\Omega (P_n(y),e^\omega_k


)\ov{(P_m(z),e^\omega_k)}\,d\mu(\omega)=(P_n(y),P_m(y)),


$$


где $k\geqslant n$, $k\geqslant m$. Значит, пространство с мерой


$\Omega$ и множество функций $\Xi=\{\wh y_\omega^n\}_{n\in\Bbb


N,\,y\in H}$ удовлетворяют условиям леммы, функция $\theta(\omega)$


из нее и будет подходящей.


\end{pf}





В дальнейшем изложении часто будем ограничиваться рассмотрением


измеримых неотрицательных обобщенных ортоподобных систем разложения,


что по доказанной теореме в случае счетно-конечного $\Omega$ не сужает


общности.





Для измеримой неотрицательной обобщенной ортоподобной системы


разложения отображение $y\longrightarrow\wh y_\omega$ по теореме


1 является унитарным (сохраняющим скалярное произведение) линейным


оператором из гильбертова пространства $H$ в пространство Лебега


$L^2(\Omega)$ и образ $H$ --- замкнутое подпространство в $L^2(\Omega)$.





\begin{defn}\hspace{-1.8mm}{\bf.}


Будем в дальнейшем при рассмотрении


измеримой неотрицательной ортоподобной системы разложения


образ $H$ в пространстве $L^2(\Omega)$ при отображении


$y\longrightarrow \wh y_\omega$ обозначать


$L^2_H(\Omega)$ или короче $L^2_H$. Это множество коэффициентов


$\wh y_\omega$ элементов $y\in H$


$$


L^2_H(\Omega)=\{c(\omega)\in L^2(\Omega): c(\omega)=\wh y_\omega


\ \; \text{для некоторого} \ \; y\in H\}.


$$


\end{defn}





Как показывают замечания 4 и 5 к определению 5, в случае систем,


соответствующих преобразованию Фурье или преобразованию Гильберта,


$L^2_H(\Omega)=L^2(\Omega)$, где $\Omega$ --- это $\Bbb R$ с $\mu$


--- обычной мерой Лебега. В общем случае $L^2_H(\Omega)$ может не


совпадать с $L^2(\Omega)$. Так конечномерные неотрицательные


ортоподобные системы в $H=\Bbb R^n$ или в $H=\Bbb C^n$ могут


состоять из $K$ векторов, где $K>n$ [2]. В этом случае изометричное


$H$ пространство $L^2_H(\Omega)$\; $n$-мерно, а содержащее его


пространство $L^2(\Omega)$\; $K$-мерно, $\Omega$ состоит из $K$


точек. Автор не знает примеров неотрицательных ортоподобных систем,


исключая случай ортонормированных систем, когда


$L^2_H(\Omega)=L^2(\Omega)$.





В случае $L^2_H(\Omega)\ne L^2(\Omega)$ естественно рассмотреть


и другие возможные по теореме 3 разложения элементов $y\in H$. При


этом возникает вопрос о существовании удовлетворяющих условиям теоремы


3 последовательностей измеримых множеств $\{\Lambda^n_k\}^\infty_{k=1}


$, что связано с измеримостью функций $c(\omega)e_n^\omega$ из формулы


(8). Введем





\begin{defn}\hspace{-1.8mm}{\bf.}


Обобщенную систему $\{\boldsymbol e^\omega


\}_{\omega\in\Omega}$ в $H$ (с системой подпространств $\{H_n\}_{n=1


}^\infty$ и индексами из $\Omega$) будем называть $L^2$-{\it измеримой


}, если для любой функции $c(\omega)$ из пространства Лебега $L^2


(\Omega)$ со значениями в $\Bbb R$ или $\Bbb C$ все функции


$c(\omega)e_n^\omega$ измеримы как функции на $\Omega$ со значениями в


$H$.


\end{defn}





\begin{thm} {\em (о $L^2$-измеримости)}.


Если $\{\boldsymbol e^\omega\}_{\omega\in\Omega}$ ---


неотрицательная обобщенная ортоподобная


система разложения в $H$, то для ее $L^2$-измеримости необходимо и


достаточно, чтобы выполнялись следующие условия$:$


\begin{itemize}


\item[а)] для каждой функции $c(\omega)\in L^2(\Omega)$, каждого


измеримого подмножества $E\subset\Omega$ конечной меры и каждого


$n\in\Bbb N$ функция $c(\omega)e_n^\omega$ почти


сепарабельнозначна на $E;$


\item[б)] для каждого $y\in H$ и каждого $n\in\Bbb N$ функция


$c(\omega)\wh y_\omega^n$ измерима на $\Omega$.


\end{itemize}


\end{thm}





Теорема непосредственно следует из ([3], с.\,166).


\smallskip





{\bf Следствие }(о $L^2$-измеримости). Если $\{\boldsymbol


e^\omega\}_{\omega\in\Omega}$ --- неотрицательная ортоподобная система


разложения в $H$, то для ее $L^2$-измеримости достаточно, чтобы она


была измерима, а $H$ сепарабельно.





\begin{lem} {\em (о существовании 


подходящих исчерпываний)}. Если


$\{\boldsymbol e^\omega\}_{\omega\in\Omega}$ --- неотрицательная


обобщенная ортоподобная $L^2$-измеримая система разложения в $H$, то


для любой функции $c(\omega)\in L^2(\Omega)$ найдется такое


исчерпывание ${\{\Omega_k}\}^\infty_{k=1}$ пространства $\Omega$, что


для каждого $n\in\Bbb N$ функция $c(\omega)e_n^\omega$ интегрируема


по Лебегу на $\Omega_ k$.


\end{lem}





\begin{pf}


Множества $\Omega_k^n=\{\omega\in


\Omega:1/k<|c(\omega)|$ и $\|c(\omega)e^\omega_n\|


<k\}\cup\{\omega\in\Omega:c(\omega)=0\}$, $\Omega_k


=\bigcap_{n=1}^\infty\Omega_k^n$ являются измеримыми, т.\,к. по


условию функции $c(\omega)$ и $c(\omega)e^\omega_n$ измеримы. Функция


$c(\omega)\in L^2(\Omega )$ отлична от нуля лишь на подмножествах


конечной меры множеств $\Omega_k^n$ и $\Omega_k$, на этих


подмножествах измеримая функция $c(\omega)e^\omega_n$ ограничена, а


значит, интегрируема.


\end{pf}





\begin{thm}\hspace{-1.8mm}{\bf.}


Если $\{\boldsymbol e^\omega\}_{\omega\in


\Omega}$ --- неотрицательная обобщенная ортоподобная $L^2$-измеримая


система разложения в $H$, то для любой функции $c(\omega)\in L^2


(\Omega)$ существует в пространстве $H$ элемент


\begin{equation}


y=\int_\Omega c(\omega)\boldsymbol e^\omega\,d\mu(\omega)


\end{equation}


{\em (см. определение 7).}


\end{thm}





Эта теорема --- непосредственное следствие теоремы 4 и леммы 3.





В дальнейшем при рассмотрении


неотрицательной обобщенной измеримой ортоподобной системы разложения


будем для любого элемента $y\in H$ множество функций $c(\omega)\in


L^2(\Omega)$, для которых выполняется равенство (10), обозначать


$L^2_y(\Omega)$ или короче $L^2_y$.





\begin{thm}\hspace{-1.8mm}{\bf.}


Пусть $\{\boldsymbol e^\omega\}_{\omega\in


\Omega}$ --- неотрицательная $L^2$-измеримая обобщенная


ортоподобная система разложения в $H$, тогда подпространство $L^2_0$


--- ортогональное дополнение $L^2_H$ в пространстве Лебега $L^2


(\Omega)$. Замкнутое подпространство $L^2_0$ сдвигом на любой элемент


множества $L^2_y$ переходит в $L^2_y$. Функция $\wh y_\omega$


является ортогональной проекцией $L^2_y$ на $L^2_H$.


\end{thm}





\begin{pf}


Из формул (1), (6) и теоремы 2 сразу следует, что $L^2_0$ ---


ортогональное дополнение $L^2_H$ и, значит, замкнутое


подпространство ([3], с.\,270). Так как разность любых


двух элементов $L^2_y$ --- элемент $L^2_0$, а сумма любого элемента


$L^2_y$ с любым элементом $L^2_0$ --- элемент $L^2_y$, то $L^2_0$


сдвигом на любой элемент $L^2_y$ переходит в $L^2_y$. Ортогональной


проекцией $L^2_y$ на $L^2_H$ будет функция из пересечения $L^2_y$ и


$L^2_H$, т.\,е. функция $\wh y_\omega$.


\end{pf}





\section{Об ортогональных проекциях ортогональных\protect\\ и


обобщенных ортоподобных систем}





Простейшим методом получения новых неотрицательных обобщенных


ортоподобных систем разложения является проектирование уже известных


систем.





\begin{thm}\hspace{-1.8mm}{\bf.}


Если $P$ --- оператор ортогонального


проектирования в гильбертовом пространстве ${\Bbb H}$ на


подпространство $H$, а $\{\boldsymbol e^\omega\}_{\omega \in\Omega}$


--- неотрицательная обобщенная ортоподобная система разложения в


$\Bbb H$ $($с системой подпространств $\{H_n\}_{n=1}^\infty$ и


индексами из $\Omega)$, то ее проекция $\{P(\boldsymbol e^\omega


)\}_{\omega\in\Omega}$, где $P(\boldsymbol e^\omega)$ ---


последовательность $\{P(e^\omega_n)\}_{n=1}^\infty$, ---


неотрицательная обобщенная ортоподобная система разложения в $H$


$($с системой подпространств $\{P(H_n)\}_{n=1}^\infty$ и индексами из


$\Omega)$ $($причем если начальная система измерима или $L^2$-измерима,


то спроектированная система соответственно также измерима или


$L^2$-измерима$)$.


\end{thm}





\begin{pf}


Так как $\{H_n\}_{n=1}^\infty$ --- система


замкнутых расширяющихся подпространств в $\Bbb H$, объединение


которых всюду плотно в $\Bbb H$, то, следовательно,


$\{P(H_n)\}_{n=1}^\infty$ --- система замкнутых расширяющихся


подпространств в $P(\Bbb H)=H$, объединение которых всюду плотно в


$H$. В [5] доказано, что проекция неотрицательной ортоподобной системы


разложения на подпространство --- неотрицательная ортоподобная система


разложения в этом подпространстве. Следовательно, $\{P(e^\omega_n)\}_{


\omega\in\Omega}$ --- неотрицательная ортоподобная система разложения


в $P(H_n)$. Так как $e^\omega_n$ --- ортогональная проекция


$e^\omega_{n+1}$ на $H_n$, то $P(e^\omega_n)$ --- ортогональная


проекция $P(e^\omega_{n+1})$ на $P(H_n)$. Согласно определениям 4--6


$\{P(\boldsymbol e^\omega)\}_{\omega\in\Omega}$ --- неотрицательная


обобщенная ортоподобная система разложения в $H$ (с системой


подпространств $\{P(H_n)\}_{n=1}^\infty$ и индексами из $\Omega$).





Если $y\in H$, то функция $\wh y_\omega$, соответствующая


начальной системе $\{\boldsymbol e^\omega\}_{\omega\in\Omega}$,


соответствует и спроектированной системе $\{P(e^\omega)\}_{\omega\in


\Omega}$ и поэтому, если начальная система измерима или


$L^2$-измерима, то спроектированная система соответственно также


измерима или $L^2$-измерима.


\end{pf}





Эта теорема позволяет из ортонормированных, ортогональных и


неотрицательных ортоподобных систем с помощью проектирования


получать примеры неотрицательных обобщенных ортоподобных систем


разложения. Естественно выяснить, когда обобщенная ортоподобная


система фактически является ортогональной системой в $H$ или проекцией


ортогональной системы из некоторого содержащего $H$ гильбертова


пространства или ортоподобной системой в $H$ (т.\,е. когда $e^\omega_n$


--- ортогональные проекции на $H_n$ ортогональной системы из $H$ или


из некоторого содержащего $H$ гильбертова пространства либо


ортогональные проекции ортоподобной системы в $H$).





\begin{thm}\hspace{-1.8mm}{\bf.}


Пусть $\{\boldsymbol e^\omega\}_{\omega\in


\Omega}$ --- неотрицательная обобщенная ортоподобная система


разложения в $H$ $($с системой подпространств $\{H_n\}_{n=1}^\infty$


и индексами из $\Omega)$. Она вся $($почти вся$)$ является некоторой


неотрицательной ортоподобной системой разложения $\{\boldsymbol


e^\omega\}\subset H$ с $\omega$, пробегающим все $($почти все$)$ $\Omega$


в том смысле, что для каждого $n\in\Bbb N$ и каждого $($почти


каждого$)$ $\omega\in\Omega$ элемент $e^\omega_n$ --- ортогональная


проекция $\boldsymbol e^\omega$ на $H_n$ тогда и только тогда, когда


для каждого $($почти каждого$)$ $\omega\in\Omega$\; $\sup\limits_{n\in\Bbb N}


\|e^\omega_n\|<+\infty$.


\end{thm}





\begin{pf}


Необходимость очевидна,


т.\,к. если для некоторого $\omega\in\Omega$ при всех $n\in\Bbb


N$ элементы $e^\omega_n$ --- ортогональные проекции $\boldsymbol


e^\omega$ на $H_n$, то $\sup\limits_{n\in\Bbb N}\|e^\omega_n\|\leqslant\|


\boldsymbol e^\omega\|<+\infty$.





Докажем достаточность. Согласно лемме 1 для каждого (почти


каждого) $\omega\in\Omega$ существует элемент $\boldsymbol e^\omega=


\lim\limits_{n\to\infty}e^\omega_n$


и $P_n(\boldsymbol e^\omega)=e^\omega_n$.


Осталось только доказать, что система $\{\boldsymbol e^\omega\}$ с


$\omega$, пробегающим все (почти все) $\Omega$, --- неотрицательная


ортоподобная система разложения в $H$. Для любого элемента $y\in H_n$


по теореме 1\; $\wh y_\omega=\lim\limits_{m\to\infty}(y,e^\omega_m


)=(y,e^\omega_n)$ и, следовательно, $\wh y_\omega=(y,\boldsymbol


e^\omega)$ для всех (почти всех) $\omega\in\Omega$.


Так как $\bigcup_{n=1}^\infty H_n$ всюду плотно в $H$, то для


любого $y\in H$\; $\wh y_\omega=(y,\boldsymbol e^\omega)$ для всех


(почти всех) $\omega\in\Omega$. По теореме 1 для любого элемента $y\in


H$ выполняется аналог равенства Парсеваля--Планшереля $\|y\|^2


=(L)\int\limits_\Omega|\wh y_\omega|^2\,d\mu(\omega)$, а по [5]


это необходимое и достаточное условие того, чтобы система


$\{\boldsymbol e^\omega\}$ была неотрицательной ортоподобной системой


разложения в $H$.


\end{pf}





\begin{lem}\hspace{-1.8mm}{\bf.}


Пусть $\{e^\omega\}_{\omega\in\Omega}$ ---


неотрицательная ортоподобная система разложения в $H$. Если $\mu


(\tau)>0$, где $\tau\in\Omega$, то $\|e^\tau\|^2\leqslant


1/{\mu(\tau)}$.


\end{lem}





\begin{pf}


По аналогу равенства Парсеваля


$$


\|e^\tau\|^2=\int_\Omega |(e^\tau,e^\omega)|^2\,d\mu(\omega


)\geqslant\|e^\tau\|^4\mu(\tau),


$$


значит, утверждение леммы верно.


\end{pf}





\begin{thm}\hspace{-1.8mm}{\bf.}


Пусть $\{\boldsymbol e^\omega\}_{\omega


\in\Omega}$ --- неотрицательная обобщенная ортоподобная система


разложения в $H$ $($с системой подпространств $\{H_n\}_{n=1}^\infty$ и


индексами из $\Omega)$. Она вся $($почти вся$)$ является некоторой полной


ортонормированной $($ортогональной$)$ системой $\{\boldsymbol e^\omega\}$


из $H$ с $\omega$, пробегающим все $($почти все$)$ $\Omega$ в том смысле,


что для каждого $n\in{\Bbb N}$ и каждого $($почти каждого$)$ $\omega\in


\Omega$ элемент $e^\omega_n$ --- ортогональная проекция $\boldsymbol


e^\omega$ на $H_n$ тогда и только тогда, когда для каждого $($почти


каждого$)$ $\omega\in\Omega$\; $\boldsymbol e^\omega=\{e^\omega_n


\}_{n\in\Bbb N}$ не является нулевой последовательностью,


$\mu(\omega)=1$ $(\mu(\omega)>0)$


и для каждого $($почти каждого$)$ $\rho\in\Omega$, для каждого $($почти


каждого$)$ $\tau\in\Omega$, $\tau\ne\rho$, $\lim\limits_{n\to\infty}


(e^\rho_n,e^\tau_n)=\int\limits_\Omega\wh{\boldsymbol e^\rho}_\omega


\ov{\wh{\boldsymbol e^\tau}_\omega}\,d\mu(\omega)=0$.


\end{thm}





\begin{pf}


Необходимость очевидна, т.\,к.





1) для ортонормированной (ортогональной) системы $\{\boldsymbol


e^\omega\}$ все индексы $\omega$ имеют меру $\mu(\omega)=1$ ($>0$);





2) скалярное произведение разных ее членов $(\boldsymbol


e^\rho,\boldsymbol e^\tau)=\lim\limits_{n\to\infty}(e^\rho_n,e^\tau_n


)=\int\limits_\Omega\wh{\boldsymbol e^\rho}_\omega\ov


{\wh{\boldsymbol e^\tau}_\omega}\,d\mu(\omega)=0$.





Докажем достаточность. Из леммы 2 и теоремы 9 следует, что





3) подсистема $\{\boldsymbol e^\omega\}_{\mu(\omega)>0}$


является некоторой неотрицательной ортоподобной системой разложения


$\{\boldsymbol e^\omega\}\subset H$ с такими $\omega\in\Omega$, что


$\mu(\omega)>0$.





Скалярное произведение двух разных ее членов $(


\boldsymbol e^\rho,\boldsymbol e^\tau){=}\lim\limits_{n\to\infty}


(e^\rho_n,e^\tau_n)=\int\limits_\Omega\wh{\boldsymbol e^\rho}_\omega


\ov{\wh{\boldsymbol e^\tau}_\omega}\,d\mu(\omega){=}0$,


значит, это ортогональная система. При этом по аналогу равенства


Парсеваля


$\|\boldsymbol e^\rho\|^2=\linebreak|(\boldsymbol e^\rho,\boldsymbol


e^\rho)|^2\mu(\rho)$ и, значит, если $\mu(\rho)=1$, то $\|\boldsymbol


e^\rho\|=1$.


\end{pf}





\begin{thm}\hspace{-1.8mm}{\bf.}


Пусть $\{\boldsymbol e^\omega\}_{\omega


\in\Omega}$ --- неотрицательная обобщенная ортоподобная система


разложения в $H$ $($с системой подпространств $\{H_n\}_{n=1}^\infty$


и индексами из $\Omega)$. Она вся $($почти вся$)$ является ортогональной


проекцией на $H$ некоторой полной ортонормированной $($ортогональной$)$


системы $\{\boldsymbol\epsilon^\omega\}$ с $\omega$, пробегающим все


$($почти все$)$ $\Omega$, из какого-то содержащего $H$ гильбертова


пространства в том смысле, что для каждого $n\in{\Bbb N}$ и каждого


$($почти каждого$)$ $\omega\in\Omega$ элемент $e^\omega_n$ ---


ортогональная проекция $\boldsymbol\epsilon^\omega$ на $H_n$ тогда и


только тогда, когда для каждого $($почти каждого$)$


$\omega\in\Omega$\; $\mu(\omega)=1$ $(\mu(\omega)>0)$.


\end{thm}





\begin{pf}


Необходимость следует из п.\,1) предыдущего доказательства.





Достаточность. Согласно п.\,3) предыдущего доказательства


и [5] система $\{\boldsymbol e^\omega\}_{\mu(\omega)>0}$ является


ортогональной проекцией на $H$ некоторой полной ортонормированной


(ортогональной) системы $\{\boldsymbol\epsilon^\omega\}_{\mu


(\omega)>0}$ из какого-то содержащего $H$ гильбертова пространства.


\end{pf}





{\bf Следствие.} Любая не более чем счетная (т.\,е. с не


более чем счетным $\Omega$) обобщенная ортоподобная система почти вся


(а если $\mu(\omega)>0$ для каждого $\omega\in\Omega$, то вся)


является ортогональной проекцией на $H$ некоторой полной ортогональной


системы из некоторого содержащего $H$ гильбертова пространства.





\section{Критерий неотрицательных обобщенных ортоподобных


систем\protect\\ --- аналог равенства Парсеваля--Планшереля}





\begin{thm}\hspace{-1.8mm}{\bf.}


Неотрицательная $L^2$-измеримая обобщенная


система $\{\boldsymbol e^\omega\}_{\omega\in\Omega}$ в $H$ $($с системой


подпространств $\{H_n\}_{n=1}^\infty$ и индексами из $\Omega)$


будет неотрицательной обобщенной ортоподобной системой разложения


в $H$ тогда и только тогда, когда для любого элемента $y\in H$


выполняется аналог равенства Парсеваля--Планшереля


$$


\|y\|^2=\lim_{n\to\infty}\int_\Omega|(y,e^\omega_n)|^2


\,d\mu(\omega).


$$


\end{thm}





\begin{pf}


Необходимость выполнения аналога равенства


Парсеваля--Планшереля видна из теоремы 1.





Докажем


достаточность. Если $y\in H_n$, то для $m>n$\ $(y,e^\omega_m)=(y,


e^\omega_n)$ и, значит, $\|y\|^2=\int\limits_\Omega|(y,e^\omega_n)|^2\,d\mu


(\omega)$. Тогда по теореме 10 из [5] $\{e^\omega_n\}_{\omega\in


\Omega}$ --- неотрицательная ортоподобная система разложения в $H_n$,


а значит, $\{\boldsymbol e^\omega\}_{\omega\in\Omega}$ ---


неотрицательная обобщенная ортоподобная система разложения в $H$.


\end{pf}





В случае счетно-конечного $\Omega$ и сепарабельного $H$


требование $L^2$-измеримости можно отбросить.





\begin{thm}\hspace{-1.8mm}{\bf.}


Неотрицательная обобщенная система


$\{\boldsymbol e^\omega\}_{\omega\in\Omega}$ в сепарабельном


гильбертовом пространстве $H$ с системой подпространств


$\{H_n\}_{n=1}^\infty$ и индексами из счетно-конечного пространства


с мерой $\Omega$ будет неотрицательной обобщенной ортоподобной


системой разложения в $H$ тогда и только тогда, когда для любого


элемента $y\in H$ выполняется аналог равенства


Парсеваля--Планшереля


$$


\|y\|^2=\lim_{n\to\infty}\int_\Omega|(y,e^\omega_n)|^2\,d\mu(\omega).


$$


\end{thm}





Доказательство повторяет доказательство теоремы 13 с заменой


ссылки на теорему 10 из [5] ссылкой на теорему 11 из [5].





\section{О конструкциях обобщенных ортоподобных систем}





Докажем ряд теорем о методах построения неотрицательных


обобщенных ортоподобных систем разложения. Простейшей из них


является





\begin{thm}\hspace{-1.8mm}{\bf.}


Если $\boldsymbol P$ --- оператор


проектирования в гильбертовом пространстве $H$ на


подпространство $H'$, а $\{\boldsymbol e^\omega\}_{\omega\in\Omega}$


--- неотрицательная обобщенная ортоподобная система разложения в $H$


$($с системой подпространств $\{H_n\}_{n=1}^\infty$ и индексами из


$\Omega)$ $($по определениям $4$--$6)$, $\boldsymbol e^\omega=\{e^\omega_n


\}_{n\in\Bbb N}$, то ее проекция $\{\boldsymbol P(\boldsymbol


e^\omega)\}_{\omega\in\Omega}$, где $\boldsymbol P(\boldsymbol


e^\omega)=\{\boldsymbol P(e^\omega_n)\}_{n\in\Bbb


N}$, --- неотрицательная обобщенная ортоподобная система разложения в


$H'$ $($с системой подпространств $\{H'_n\}_{n=1}^\infty$, где $H'_n


=\boldsymbol P(H_n)$, и индексами из $\Omega)$.


\end{thm}





\begin{pf}


Если $P_n$ --- оператор ортогонального проектирования из $H$ на $H_n$,


то композиция $\boldsymbol PP_n$ является оператором ортогонального


проектирования из $H'$ на $H'_n$ и $\boldsymbol PP_n(e^\omega_{


n+1})=\boldsymbol P(e^\omega_n)$, а согласно [5] $\{\boldsymbol


P(e^\omega_n)\}_{\omega\in\Omega}$ --- неотрицательная ортоподобная


система разложения в $H'_n$.


\end{pf}





Укажем еще один метод построения неотрицательных обобщенных


ортоподобных систем разложения.





\begin{thm}\hspace{-1.8mm}{\bf.}


Пусть $\{\boldsymbol e^\omega\}_{\omega\in


\Omega}$ --- неотрицательная обобщенная ортоподобная система


разложения в $H$ $($с системой подпространств $\{H_n\}_{n=1}^\infty$


и индексами из $\Omega)$ $($по определениям $4$ и $5)$, $\boldsymbol


e^\omega=\{e^\omega_n\}_{n\in\Bbb N}$. Пусть $\Upsilon$ ---


пространство со счетно-аддитивной неотрицательной мерой $\nu$,


$\Omega\times\Upsilon$ --- произведение пространств $\Omega$ и


$\Upsilon$, $\mu\times\nu$ --- произведение мер $\mu$ и $\nu$


{\em ([3], с.\,201--206)}, $\Phi_\upsilon$ ---


такое измеримое отображение $\Upsilon$ в множество линейных


непрерывных операторов, отображающих $H$ в $H$ и коммутирующих с


операторами $P_n$ ортогонального проектирования $H$ на $H_n$, что


\begin{equation}


\int_{\Upsilon}\Phi_\upsilon\cdot\Phi_\upsilon^*\,d\nu(\upsilon)=E,


\end{equation}


$E$ --- тождественный оператор, под $\Phi^*$ понимается оператор,


сопряженный с $\Phi$, интеграл берется от функций со значениями


в нормированном пространстве линейных непрерывных операторов. Тогда


система $\{\Phi_\upsilon(\boldsymbol e^\omega)\}_{(\omega


,\upsilon)\in\Omega\times\Upsilon}$, где $\Phi_\upsilon(\boldsymbol


e^\omega)=\{\Phi_\upsilon(e^\omega_n)\}_{n\in\Bbb N}$, ---


ортоподобная неотрицательная система разложения в $H$


$($с системой подпространств $\{H_n\}_{n=1}^\infty$ и индексами из


$\Omega\times\Upsilon)$.


\end{thm}





\begin{pf}


Сначала покажем, что для каждого


натурального $n$\; $\{\Phi_\upsilon


(e^\omega_n)\}_{(\omega,\upsilon)\in\Omega\times\Upsilon}$ ---


неотрицательная ортоподобная система разложения в $H_n$.





Пусть $y\in H_n$, т.\,к. $(y,\Phi_\upsilon (e^\omega_n


))=(\Phi^*_\upsilon(y),e^\omega_n)$, то


$$


\Phi^*_\upsilon (y)=\int_\Omega (y,\Phi_\upsilon


(e^\omega_n))e^\omega\,d\mu(\omega).


$$


Так как $\Phi_\upsilon$ --- непрерывный линейный оператор, то,


используя ([3], с.\,128), имеем


$$


\Phi_\upsilon\cdot\Phi^*_\upsilon (y)=\int_\Omega(y,\Phi_\upsilon


(e^\omega_n))\Phi_\upsilon (e^\omega_n)\,d\mu(\omega)


$$


как для собственного, так и несобственного интегрирования по Лебегу.


В силу (11)


$$


y=(L)\int_\Upsilon \Phi_\upsilon\cdot\Phi^*_\upsilon (y)\,d\nu


(\upsilon)=(L)\int_\Upsilon \bigg(\int_\Omega(y,\Phi_\upsilon


(e^\omega_n))\Phi_\upsilon (e^\omega_n)\,d\mu(\omega)\bigg)d\nu


(\upsilon).


$$


Теорема доказана в случае, когда внутренний интеграл в правой части


последнего равенства --- собственный интеграл Лебега. В случае, когда


этот интеграл несобственный, требуется дополнительное рассуждение.


По условию теоремы $\Phi_\upsilon\cdot \Phi^*_\upsilon$, а значит, и


$\|\Phi_\upsilon \cdot \Phi^*_\upsilon\|$


интегрируема по Лебегу на $\Upsilon$. По утверждениям 1 и 3


$$


\bigg\|\int_{\Omega_k} (y, \Phi_\upsilon (e^\omega_n))e^\omega_n\,


d\mu(\omega)\bigg\|=\bigg\|\int_{\Omega_k}(\Phi^*_\upsilon


(y),e^\omega_n)e^\omega_n\,d\mu(\omega)\bigg\|\leqslant\|


\Phi^*_\upsilon(y)\|\leqslant \|\Phi^*_\upsilon\| \cdot \|y\|.


$$


Так как для любого линейного непрерывного оператора $A$ из $H$ в $H$


выполняется равенство


$$


\|A\cdot A^*\|=\sup_{\|y\|\leqslant 1,\ \|z\|\leqslant 1}


(A\cdot A^*(y),z)=\sup_{\|y\|\leqslant 1,\ \|z\|\leqslant


1} (A^*(y),A^* (z))=\|A^*\|^2=\|A\|^2,


$$


то значит,


$$


\bigg\|\Phi_\upsilon\bigg(\int_{\Omega_k}(y,\Phi_\upsilon


(e^\omega_n))e^\omega_n\,d\mu(\omega)\bigg)\bigg\|\leqslant\|


\Phi_\upsilon\|\cdot\|\Phi^*_\upsilon\|\cdot\|y\|=\|\Phi_\upsilon


\cdot\Phi^*_\upsilon\|\cdot\|y\|,


$$


справа стоит интегрируемая по Лебегу на $\Upsilon$ функция. Пользуясь


теоремой Лебега о предельном переходе под знаком интеграла ([3],


с.\,168), получаем


\begin{gather*}


y=(L)\int_\Upsilon \Phi_\upsilon\cdot\Phi^*_\upsilon (y)\,


d\nu(\upsilon)=(L)\int_\Upsilon \bigg(\lim_{k\to\infty} \int


_{\Omega_k}(y,\Phi_\upsilon(e^\omega_n))\Phi_\upsilon


(e^\omega_n)\,d\mu(\omega)\bigg)d\nu(\upsilon)= \\


%


=\lim_{k\to\infty}(L)\int_\Upsilon \int_{\Omega_k}


(y,\Phi_\upsilon (e^\omega_n))\Phi_\upsilon(e^\omega_n)\,d\mu(\omega


)d\nu(\upsilon)= \\


%


=\lim_{k\to\infty} (L)\int_{\Upsilon\times\Omega_k}


(y,\Phi_\upsilon(e^\omega_n))\Phi_\upsilon(e^\omega_n)\,d(\nu(\upsilon


)\times\mu(\omega)).


\end{gather*}


Следовательно, для каждого натурального $n$\; $\{\Phi_\upsilon


(e^\omega_n)\}_{(\omega,\upsilon)\in\Omega\times\Upsilon}$ ---


неотрицательная ортоподобная система разложения в $H_n$.





По условию операторы $\Phi_\upsilon$ и $P_n$ коммутируют,


поэтому $P_n(\Phi_\upsilon(e^\omega_{n+1}))=\Phi_\upsilon


(P_n(e^\omega_{n+1}))=\Phi_\upsilon(e^\omega_n)$.





По определениям 4--6 система $\{\Phi_\upsilon(\boldsymbol


e^\omega)\}_{(\omega,\upsilon)\in\Omega\times\Upsilon}$,


где $\Phi_\upsilon(\boldsymbol e^\omega)=\{\Phi_\upsilon(e^\omega_n)


\}_{n\in\Bbb N}$, ---


ортоподобная неотрицательная система разложения в $H$


(с системой подпространств $\{H_n\}_{n=1}^\infty$ и индексами из


$\Omega\times\Upsilon$).


\end{pf}





Наиболее известными континуальными неотрицательными ортоподобными


системами разложения являются всплески Д.\,Габора и всплески Дж.\,Морле.


Покажем, что они строятся по доказанной теореме на основе


неотрицательной обобщенной ортоподобной системы функций


$\{\exp(2\pi i\omega x)\}_{\omega\in\Bbb R}$,


определяющей преобразование Фурье (см. замечание 4 к определению 6).


\medskip





{\it Всплески Габора.} Порожденное одной функцией $w(x)$


семейство функций


$$


w_{\lambda, \beta}(x)=e^{2\pi i\lambda x} w(x-\beta),


$$


где параметры $\lambda$, $\beta$ и переменная $x$ из ${\Bbb R}$,


называют системой всплесков Габора. Если функция $w\in L^2 ({\Bbb


R})$ и $\|w\|_{L^2 ({\Bbb R})}=1$, то для любой функции $f(x)\in


L^2 ({\Bbb R})$


$$


f(x)=\int_{\Bbb R}\int_{\Bbb R} (f, w_{\lambda,


\beta})w_{\lambda, \beta}(x)\,d\lambda\,d\beta,


$$


где последний интеграл понимается как предел в $L^2(\Bbb R)$


$$


\int_{\Bbb R}\int_{|\lambda|\leqslant A} (f, w_{\lambda,


\beta})w_{\lambda, \beta}(x)\,d\lambda\,d\beta


$$


при $A\to\infty$.





Этот результат при дополнительном предположении суммируемости


на ${\Bbb R}$ функции $w(x)$ был получен Д.\,Габором [9] (или см. [10],


с.\,317--320). В более общем виде он изложен в [11], [12] и в [6].





Если для $\beta\in\Bbb R$ взять операторы


$\Phi_\beta(f)(x)=f(x)w(x-\beta)$, то


$$


\int_{\Bbb R}\Phi_\beta\cdot\Phi_\beta^*\,d\beta=\int_{\Bbb


R}|w(\beta)|^2\,d\beta\cdot E=E,


$$


где $E$ --- тождественный оператор, интеграл берется от функций со


значениями в нормированном пространстве линейных непрерывных


операторов. Используя теоремы 14 и 9, можно констатировать, что


семейство всплесков Габора


$\{w_{\lambda,\beta}(x)\}_{(\lambda,\beta)\in\Bbb R^2}$


с обычной мерой Лебега на $\Bbb R^2$ будет неотрицательной


ортоподобной системой разложения.


\medskip





{\it Всплески Морле.} Порожденное одной функцией $\psi(x)$


семейство функций


$$


\psi_{\alpha, \beta}(x)=\frac1{\sqrt{|\alpha|}}\psi\bigg(\frac


{x-\beta}{\alpha}\bigg),


$$


где параметр $\alpha\in{\Bbb R}\setminus\{0\}$, параметр $\beta\in


{\Bbb R}$ и переменная $x\in{\Bbb R}$, называют системой


всплесков Морле. Если функция $\psi\in L^2(\Bbb R)$ и


$$


\int_{\Bbb R}\frac{|\wh\psi(\lambda)|^2}{|\lambda|}\,


d\lambda=1,


$$


где $\wh\psi(\lambda)$ --- преобразование Фурье функции $\psi(x)$,


то для любой функции $f(x)\in L^2(\Bbb R)$


$$


f(x)=\int_{\Bbb R}\int_{\Bbb R} (f, \psi_{\alpha,


\beta})\psi_{\alpha, \beta}(x)\,d\beta\,\frac{d\alpha}{\alpha^2},


$$


где последний интеграл понимается как предел в $L^2(\Bbb R)$


$$


\int_{|\alpha|\geqslant\varepsilon}\int_{\Bbb R} (f,\psi_{\alpha,


\beta})\psi_{\alpha, \beta}(x)\,d\beta\,\frac{d\alpha}{\alpha^2}


$$


при $\varepsilon\to+0$.





Этот результат при дополнительном предположении суммируемости


на ${\Bbb R}$ функции $\psi(x)$ был получен А.\,Гроссманном и Дж.\,Морле


[13] (или [10], с.\,325--329). В более общем виде он изложен в


[11], [12] и в [6].





Если взять для $\alpha\in\Bbb R\setminus\{0\}$ операторы


$\Phi_\alpha(f)(x)={\sqrt{|\alpha|}}\,\wh\psi(\alpha x)f(x)$, то


$$


\int_{\Bbb R}\Phi_\alpha\cdot\Phi_\alpha^*\,


\frac{d\alpha}{\alpha^2}=\int_{\Bbb R}


|\wh\psi(\alpha x)|^2\,\frac{d\alpha}{\alpha}\cdot E=E,


$$


где $E$ --- тождественный оператор, интеграл берется от функций со


значениями в нормированном пространстве линейных непрерывных


операторов, а $\wh\psi$ ---


преобразование Фурье функции $\psi(x)$. Используя теоремы 14


и 9, можно констатировать, что семейство всплесков


$\{\sqrt{|\alpha|}\,\wh\psi(\alpha x)\text{exp}\,2\pi i\beta x\,


\}_{(\alpha,\beta)\in\Bbb R^2}$ с мерой на $\Bbb R^2$, задаваемой


дифференциалом $\frac{d\alpha}{\alpha^2}d\beta$,


будет неотрицательной ортоподобной системой разложения.


Применив к ней унитарный (сохраняющий скалярное произведение) оператор


--- преобразование Фурье, получим неотрицательную ортоподобную систему


$\{\psi_{\alpha,\beta}(x)\}_{(\alpha,\beta)\in(\Bbb R\setminus\{0\}


)\times\Bbb R}$ --- систему всплесков Морле.





Заметим, что в [6] оба приведенных примера также получаются при


помощи некоторой общей конструкции, связанной с унитарным оператором


--- преобразованием Фурье.





В заключение статьи отметим, что ряд построенных в [6] систем


разложения можно получить аналогичным образом, исходя из


неотрицательных обобщенных ортоподобных систем, связанных с


преобразованием Гильберта (см. замечание 5 к определению 4) или


с преобразованием Фурье на локально-компактной топологической группе,


или с преобразованием Рисса.





\begin{thebibliography}{99}





\bibitem{1}


Лукашенко Т.П. {\em О системах разложения, подобных


ортогональным} // Международн. конф. по теор. приближ. функций, посв.


памяти проф. П.П.\,Коровкина. Калуга, 26--29 июня 1996~г. Тез. докл.


Т.\,2. -- Тверь: Изд--во Тверск. ун-та, 1996. -- С.\,135--136.


\bibitem{2}


Лукашенко Т.П. {\em Системы разложения, подобные


ортогональным} // Фундамент. и прикл. матем. -- 1997. --


Т.\,3. -- \No\,2. -- С.\,487--517.


\bibitem{3}


Данфорд Н., Шварц Дж.Т. {\em Линейные операторы.


Общая теория}. Ч.\,1. -- М.: Ин. лит., 1962. -- 895~с.


\bibitem{4}


Лукашенко Т.П. {\em Ортоподобные неотрицательные


системы разложения} // Вестн. Моск. ун-та. Сер. матем., мех. -- 1997.


-- \No\,5. -- С.\,27--31.


\bibitem{5}


Лукашенко Т.П. {\em О коэффициентах систем


разложения, подобных ортогональным} // Матем. сб. -- 1997. -- Т.\,188.


-- \No\,12. -- С.\,57--72.


\bibitem{6}


Лукашенко Т.П. {\em Системы разложения на


пространствах с мерой} // Изв. РАН. Сер. матем. -- 1996. -- Т.\,60.


-- \No 1. -- С.\,165--174.


\bibitem{7}


Стейн И., Вейс Г. {\em Введение в гармонический


анализ на евклидовых пространствах}. -- М.: Мир, 1974. -- 336~с.


\bibitem{8}


Эдвардс Р. {\em Функциональный анализ}. -- М.: Мир,


1989. -- 1071~с.


\bibitem{9}


Gabor D. {\em Theory of communication} // J. Inst.


Elec. Eng. -- London, 1946. -- V.\,93. -- \No\,3. -- P.\,429--457.


\bibitem{10}


Gasquet C., Witomski P. {\em Analyse de Fourier


et applications}. -- Paris: Masson, 1990. -- 354~p.


\bibitem{11}


Лукашенко Т.П. {\em Всплески на топологических


группах} // Докл. РАН. -- 1993. -- Т.\,332. -- \No\,1. -- С.\,15--17.


\bibitem{12}


Лукашенко Т.П. {\em Всплески на топологических группах} // 


Изв. РАН. Сер. матем. -- 1994. -- Т.\,58. -- \No\,3. -- С.\,88--102.


\bibitem{13}


Grossmann A., Morlet J. {\em Decomposition of Hardy


functions into square integrable wavelets of constant shape} // SIAM J.


Math. Anal. -- 1984. -- V.\,15. -- \No\,4. -- P.\,723--736.





\end{thebibliography}





\vskip 2cm





\postup{\it Московский государственный}{$\qquad$\it Поступили}


\postup{\it университет им. М.В.\,Ломоносова}%


{\it первый вариант $30.11.1998$}


\postup{}{\it окончательный вариант $11.05.2000$}


\label{end}


\end{document}





